\documentclass{beamer} %Remove handout to reenable pauses
\usetheme{Berlin}
\usepackage{array}
\usepackage{color}
\usepackage{graphicx}
\usepackage{tikz}
\usepackage{tikz-cd}
\usepackage{wrapfig}
\usepackage{amsmath}
\usepackage{multicol}

\usetikzlibrary{matrix, shapes, trees, arrows, decorations.pathmorphing, fit}
\usepackage[utf8]{inputenc}

\DeclareMathOperator{\K}{K}
\DeclareMathOperator{\G}{G}
\DeclareMathOperator{\SL}{SL}
\DeclareMathOperator{\Sp}{Sp}
\DeclareMathOperator{\GL}{GL}
\DeclareMathOperator{\E}{E}
\DeclareMathOperator{\St}{St}
\DeclareMathOperator{\Cent}{Cent}
\DeclareMathOperator{\sr}{sr}
\DeclareMathOperator{\rk}{rk}
\DeclareMathOperator{\Img}{Im}
\DeclareMathOperator{\Umd}{Umd}
\DeclareMathOperator{\Ker}{Ker}
\DeclareMathOperator{\coker}{coker}
\newcommand{\catname}[1]{{\normalfont\textbf{#1}}}
\newcommand{\Groups}{\catname{Groups}}
\newcommand{\rA}{\mathsf{A}}
\newcommand{\rB}{\mathsf{B}}
\newcommand{\rC}{\mathsf{C}}
\newcommand{\rD}{\mathsf{D}}
\newcommand{\rE}{\mathsf{E}}
\newcommand{\rF}{\mathsf{F}}
\newcommand{\rG}{\mathsf{G}}
\newcommand{\StH}{\widehat{H}}
\DeclareMathOperator{\UU}{U}

\theoremstyle{definition}
\newtheorem*{thm}{Theorem}
\newtheorem*{prop}{Proposition}
\newtheorem*{lem}{Lemma}
\newtheorem*{exm}{Example}
\newtheorem*{dfn}{Definition}
\newtheorem*{que}{Question}
\newtheorem*{conj}{Conjecture}
\newtheorem*{ans}{Answer}
\newtheorem*{rem}{Remark}
\newtheorem*{cor}{Corollary}
\newtheorem*{conc}{Conclusion}
\newtheorem*{fac}{Fact}
\newtheorem*{ide}{Idea}
\newtheorem*{prb}{Problem}

\def\hnode{*[o]{\circ}}
\def\lnode#1{*[o]{\composite{!{\circ}*h!<0pt,-1em>{\scriptstyle #1}}}}
\def\lnoded#1{*[o]{\composite{!{\circ}*h!<0pt,0.5em>{\scriptstyle #1}}}}
\def\lbnode#1{*[o]{\composite{!{\bullet}*h!<0pt,-1em>{\scriptstyle #1}}}}
\def\harr#1{\ar@{-}[r]_{\displaystyle #1}}
\def\hard{\ar@{-}[d]}
\def\harl#1{\ar@{-}[l]^{\displaystyle #1}}

\newcommand\restr[2]{{\left.\kern-\nulldelimiterspace #1\vphantom{\big|}\right|_{#2}}}

\tikzset{
  root/.style={
    circle,
    draw,
    fill=black,
    inner sep=1.5pt
  },
  highlighted/.style={
    fill=green!60!black
  },
  zeroroot/.style={
    fill=white
  },
  levi/.style={
    draw,
    rounded corners,
    inner sep=3pt
  },
  dottededge/.style={
    dotted,
    thick
  }
}

\title{Horrocks theorem for \texorpdfstring{K2}{$К_2$}, weight elements in Chevalley groups and Curtis--Tits decompositions}
\author{Sergei Sinchuk}
\date{3 April 2026}

\AtBeginSection[]{}

%%%%%%%%%%%%%%%%%%%%%
\begin{document}

\begin{frame}
\titlepage
\end{frame}

\section{Recollection}

\begin{frame} \tableofcontents[currentsection] \end{frame}

\begin{frame}{Recollection}
Let $A$ be a local comm., unital domain.
\begin{prop}[Suslin, 1977]
 \( \E(n, A[X^{\pm 1}]) = \E(n, A[X]) \cdot B(A[X^{\pm 1}]) \cdot \E(n, A[X^{\pm 1}], M[X^{\pm 1}]). \)
\end{prop}
\pause
\textbf{Idea:} Denote by $L$ the Minkowski-product of three groups in the right-hand side.
Set $\delta_i(X) := \mathrm{diag}(1, \ldots, \underbrace{X}_{i\text{-th}}, \ldots, 1)$, $1 \leq i \leq n$.

\pause Show that $\delta_i^{-1} L \delta_i \subseteq L$.

\pause It is even enough to show that $\delta_i^{-1} \E(n, A[X]) \delta_i \subseteq L$.
\end{frame}

\begin{frame}[fragile]{Recollection}
\begin{columns}[T]
\begin{column}{0.4\textwidth}
\raggedleft
\scriptsize
\hspace*{0pt}%
\[
\begin{tikzcd}[row sep=2.6em, column sep=3em]
G_{123}
  \arrow[rr, "f_{23}"]
  \arrow[dd, "f_{12}"']
  \arrow[rd, "f_{13}"]
&&
G_{23}
  \arrow[dd, dashed, "{g_2^3}" description, pos=0.3]
  \arrow[rd, "{g_3^2}", pos=0.3]
& \\
&
G_{13}
  \arrow[rr, "{g_3^1}" description, pos=0.3]
  \arrow[dd, "{g_1^3}" description, pos=0.3]
&&
G_3
  \arrow[dd, "h_3"]
\\
G_{12}
  \arrow[rr, dashed, "{g_2^1}" description, pos=0.3]
  \arrow[rd, "{g_1^2}", pos=0.3]
&&
G_2
  \arrow[rd, dashed, "h_2"']
& \\
&
G_1
  \arrow[rr, "h_1"']
&&
G
\end{tikzcd}
\]
\end{column}

\begin{column}{0.6\textwidth}
\raggedleft
\scriptsize
\[
\begin{aligned}
V \;&=\; G_1 \times G_2 \times G_3,\\
W \;&=\; G_{12} \times G_{13} \times G_{23},\\
W \;&\curvearrowright \; V\text{ via }\star\text{-action},\\
\star \;&\; \text{defines an equiv. relation $\sim_W$},\\
V/\sim_W\; &\rightarrow \; G\text{ induced by }h_1, h_2, h_3,\\
V/\sim_W\;&\;\text{ is a set with no add. structure, }\\
\text{class of} \;&\;(a, b, c)\text{ is denoted }[a, b, c].\\
\end{aligned}
\]
\end{column}
\end{columns}

\scriptsize { \[(a, b, c) \star (x, y, z) = (x \cdot g_1^3(b) \cdot g_1^2(a), g_2^1(a)^{-1} \cdot y \cdot g_2^3(c),\ g_3^2(c)^{-1} \cdot g_3^1(b)^{-h_2(y)} \cdot z).\] }


\end{frame}

\begin{frame}[fragile]{Recollection}
\begin{columns}[T]
\begin{column}{0.4\textwidth}
\raggedleft
\scriptsize
\hspace*{0pt}%
\[
\begin{tikzcd}[row sep=1.8em, column sep=2.1em]
U^+_M
  \arrow[rr, hook]
  \arrow[dd, hook]
  \arrow[rd, hook]
&&
B_M
  \arrow[dd, dashed]
  \arrow[rd, hook]
& \\
&
G_M^+
  \arrow[rr]
  \arrow[dd, "\mu", pos=0.3]
&&
G_M
  \arrow[dd, "\mu"]
\\
U^+
  \arrow[rr, hook, dashed]
  \arrow[rd, hook]
&&
B
  \arrow[rd, hook, dashed]
& \\
&
G^+
  \arrow[rr]
&&
G
\end{tikzcd}
\]
\end{column}

\begin{column}{0.6\textwidth}
\raggedleft
\scriptsize
\[
\begin{aligned}
G \;&=\; \St(\Phi, A[X,X^{-1}]),\\
G^+ \;&=\; \St(\Phi, A[X]),\\
B \;&=\; \UU(\Phi^+, A[X,X^{-1}])
       \rtimes \StH(\Phi, A[X,X^{-1}]) \le G,\\
G_M \;&=\; \St(\Phi, A[X,X^{-1}], M[X,X^{-1}]),\\
U^+ \;&=\; \UU(\Phi^+, A[X]),\\
G_M^+ \;&=\; \St(\Phi, A[X], M[X]),\\
B_M \;&=\; \UU(\Phi^+, M[X,X^{-1}]) \times \{X,1+M\} \le G_M,\\
U_M^+ \;&=\; \UU(\Phi^+, M[X]).
\end{aligned}
\]
\end{column}
\end{columns}
\medbreak

\scriptsize { Here $\StH(\Phi, A[X^{\pm 1}]) = \langle h_\alpha(u) \mid u \in A[X^{\pm 1}]^\times, \alpha \in \Phi \rangle \leq \St(\Phi, A[X, X^{-1}])$. }
\pause

\scriptsize { We have $V = G^+ \times B \times G_M$, $W = U^+ \times G_M^+ \times B_M^+$. }
\end{frame}

\begin{frame}{The main result}
\begin{thm}
 The group $G = \St(\Phi, A[X, X^{-1}])$ acts on $V/\sim_W$, this action is compatible with the action of $G^+$ and, moreover, for every $z \in G_M$ one has $\mu(z) \cdot [1, 1, 1] = [1, 1, z]$, where $\mu \colon G_M \to G$.
\end{thm}
\pause The particular case $\Phi=\rA_{\geq 4}$ is Tulenbaev's Proposition 4.1.

\pause \textbf{Idea:} $G^+$ acts on the first coordinate of $V$, this action is compatible with $\star$ and descends to $V/\sim_W$.

\pause It remains to define a conj. action of $\delta_i$ (or their analogue) on $V$ and then express $G$ in terms of $G^+$ and $\delta_i$.
\end{frame}

\begin{frame}{Abe's automorphism $\sigma$}
Let $\Phi= \mathsf{D}_\ell, \mathsf{E}_6$ or $\mathsf{E}_7$.
Let $k$ be the green vertex on the diagram:

\medbreak
\begin{tabular}{ccc}
\scalebox{0.52}{%
\begin{tikzpicture}
\node[root, highlighted, label=$1$] (d1) at (0,0.5) {};
\node[root, label=$2$] (d2) at (1,0) {};
\node[root, label=$3$] (d3) at (2,0) {};
\node at (3,0) {\ldots};
\node[root, label={[xshift=-0.5cm,yshift=-0.6cm]$\ell-2$}] (dl2) at (4,0) {};
\node[root, label={[xshift=-0.6cm,yshift=-0.35cm]$\ell-1$}] (dl1) at (5,0.5) {};
\node[root, label={[xshift=-0.5cm,yshift=-0.3cm]$\ell$}] (dl) at (5,-0.5) {};
\node[root, zeroroot, label=below:{$0$}] (d0) at (0,-0.5) {};

\draw (d1) -- (d2) -- (d3) -- (2.5,0);
\draw (3.5,0) -- (dl2) -- (dl1);
\draw (dl2) -- (dl);
\draw[dottededge] (d2) -- (d0);

\node[levi, fit=(d2) (d3) (dl) (dl1) (dl2), label=above:{$\Delta$}] {};
\end{tikzpicture}}
&
\scalebox{0.52}{%
\begin{tikzpicture}
\node[root, highlighted, label=$1$] (e1) at (0,0) {};
\node[root, label=$3$] (e3) at (1,0) {};
\node[root, label={[xshift=-0.3cm]$4$}] (e4) at (2,0) {};
\node[root, label=$5$] (e5) at (3,0) {};
\node[root, label=$6$] (e6) at (4,0) {};
\node[root, label={[xshift=-0.3cm,yshift=-0.3cm]$2$}] (e2) at (2,1) {};
\node[root, zeroroot, label=left:$0$] (e0) at (2,2) {};

\draw (e1) -- (e3) -- (e4) -- (e5) -- (e6);
\draw (e2) -- (e4);
\draw[dottededge] (e2) -- (e0);

\node[levi, fit=(e2) (e3) (e4) (e5) (e6), label=above:{$\Delta$}] {};
\end{tikzpicture}}
&
\scalebox{0.52}{%
\begin{tikzpicture}
\node[root, zeroroot, label={$0$}] (f0) at (0,0) {};
\node[root, label={$1$}] (f1) at (1,0) {};
\node[root, label={$3$}] (f3) at (2,0) {};
\node[root, label={[xshift=-0.3cm]$4$}] (f4) at (3,0) {};
\node[root, label={$5$}] (f5) at (4,0) {};
\node[root, label={$6$}] (f6) at (5,0) {};
\node[root, label={[xshift=-0.3cm,yshift=-0.3cm]$2$}] (f2) at (3,1) {};
\node[root, highlighted, label={$7$}] (f7) at (6,0) {};

\draw (f1) -- (f3) -- (f4) -- (f5) -- (f6) -- (f7);
\draw (f2) -- (f4);
\draw[dottededge] (f0) -- (f1);

\node[levi, fit=(f1) (f2) (f3) (f4) (f5) (f6), label=above:{$\Delta$}] {};
\end{tikzpicture}}
\\[-0.2em]
$\mathsf{D}_\ell$ & $\mathsf{E}_6$ & $\mathsf{E}_7$
\end{tabular}
\medbreak
\pause

$m_k(\alpha_\mathrm{\max}) = 1$, therefore $m_k(\alpha) = 1$ for all $\alpha \in \Sigma_k^+$ and the subgroup $\UU(\Sigma^+_k, R)$ is abelian.
\pause
Also the weight $\varpi_k$ is a \textit{microweight} of $\Phi$.

\pause E. Abe defines $\sigma (x_\alpha(u)) := x_\alpha(u)^{h_{\omega_k}(X)} = x_\alpha(X^{\langle \varpi_k, \alpha \rangle}u)$ and shows that $\sigma(L) \subseteq L$.
\end{frame}

\begin{frame}{Comparison with Tulenbaev}
 \begin{itemize}
  \pause \item Tulenbaev follows Suslin's approach, so he constructs automorphisms $\sigma_i \colon V/\sim_W \to V/\sim_W$ ($1\leq i \leq n$) ``modeling'' the diagonal automorphisms of the Steinberg group $G = \St(n, R[X, X^{-1}])$.
  \pause \item He is then forced to check that $\sigma_i$, $\sigma_j$ commute between each other;
  \item He defines $x_{ij}(aX^{-1}) := \sigma_i \cdot x_{ij}(a) \cdot \sigma_i^{-1}$, and interprets $\St(n, R[X, X^{-1}])$ indirectly in these terms;
  \pause \item This all introduces immense technical strain in his proof while also leaving many gaps (e.g. checking commutator formulae);
  \pause \item This can not be generalized to Chevalley groups.
 \end{itemize}
\end{frame}

\section{Plan of the proof}
\begin{frame} \tableofcontents[currentsection] \end{frame}

\begin{frame}{Weight automorphisms}
 To $\omega \in P(\Phi)$ and $\beta \in \mathbb{Z} \Phi$ associate the product $\langle \omega, \beta\rangle \in \mathbb{Z}$.
\pause

 For $u \in S^\times$ and $\omega \in P(\Phi)$ there exists a bijection of the generating set $\St(\Phi, S)$ via the following mapping:
 \[ x_\alpha(a) \mapsto x_\alpha(u^{\langle\omega, \alpha\rangle} \cdot a),\ \alpha\in \Phi,\ a \in S. \]
\pause This bijection preserves Chevalley commutator formulae.

\pause This gives rise to an automorphism $\chi_{\omega, u}$ of $\St(\Phi, S)$.

\pause One can construct the ``relative'' version of this automorphism $\widetilde{\chi}_{\omega, X}$ on $G_M = \St(\Phi, S, I)$ using standard technique.

\pause Now set $S := R[X, X^{-1}]$, $u := X$.
\end{frame}

\begin{frame}{Definition of an automorphism of $V/\sim_W$}
 \begin{dfn} \label{sigma-def}
Let $v = (x, y, z) \in G^+ \times B \times G_M$.
We want to define Steinberg-group level analogue of Abe's automorphism $\sigma$ on $V/\sim_W$.
We need to define it first on $V$. We need something like this:
\begin{equation} \label{eq:sigma-def} (x, y, z) \mapsto (?, ? \cdot \chi_{\omega, X}(? \cdot y),\  \widetilde{\chi}_{\omega, X}(z)), \end{equation}
where $x \in G^+, y \in B$, $z \in G_M$, and
$\chi_{\omega, X}$, $\widetilde{\chi}_{\omega, X}$ are the automorphisms defined above.
\end{dfn}
\pause \textbf{Problem: } Can't define $\sigma$ on $G^+$.
\end{frame}

\begin{frame}{What should be the domain of definition of $\sigma$ on $G^+$? }
Let $\omega \in P(\Phi)$. Consider the subset $\mathcal{X}_\omega$ of the generating set of $\St(\Phi, A[X])$ consisting of those generators $x_{\alpha}(a) \in \St(\Phi, A[X])$ for which
$\langle \omega, \alpha\rangle < 0$ implies that $a \in A[X]$ is divisible by $X^{-\langle \omega, \alpha\rangle}$.
\pause

If $\langle \omega, \alpha\rangle \geq 0$ then $\mathcal{X}_\omega$ contains $x_\alpha(a)$ for all $a \in A[X]$.
\pause

Denote by $N_{\omega}$ the subgroup of $\St(\Phi, A[X])$ generated by $\mathcal{X}_\omega$.
\pause

Let $\omega \in P(\Phi)$ be a weight.

By definition, an \textit{$\omega$-pair} is a pair of mutually inverse group homomorphisms
\( N_\omega \overset{\delta_\omega}{\underset{\delta_{-\omega}}{\rightleftarrows}} N_{-\omega} \) satisfying
\begin{equation} \label{eq:delta} \delta_{\pm\omega}(x_\alpha(f)) = x_\alpha\!\bigl(X^{\pm\langle\omega,\alpha\rangle}\cdot f\bigr), \qquad x_\alpha(f)\in \mathcal{X}_{\pm\omega}. \end{equation}
\end{frame}

\begin{frame}{The subgroups $N_\omega$}
 At most one pair of homomorphisms $\delta_{\pm \omega}$ satisfying~\eqref{eq:delta} may exist.

 How to show that such $\delta_{\pm \omega}$ exist?

 \pause Recall that $\St(\Phi, A[X]) = \St(\Phi, A[X], XA[X]) \rtimes \St(\Phi, A)$.

 \pause Now if $\omega = \pm \varpi_k$ is a microweight (i.\,e. $\langle\omega, \Phi^+ \rangle \subseteq \{0, 1\}$), then
 $N_{\pm \varpi_k} = \St(\Phi, A[X], XA[X]) \rtimes P_k$, where $P_k^{\pm}$ is Steinberg parabolic subgroup corresponding to the $k$-th root.

 \pause We may use ``another presentation'' in the construction of $\delta_{\pm \omega}$.
\end{frame}


\begin{frame} [fragile]{The $\delta_{\omega}$-diagram}
\[
\begin{tikzcd}[ampersand replacement=\&, column sep=3.2em, row sep=2.6em]
N_\omega
  \arrow[r, "\delta_\omega"]
  \arrow[d, hook]
  \arrow[rr, bend left=30, "\mathrm{id}"]
\&
N_{-\omega}
  \arrow[r, "\delta_{-\omega}"]
  \arrow[d, hook]
\&
N_\omega
  \arrow[d, hook]
\\
\St(\Phi,A[X])
  \arrow[d]
\&
\St(\Phi,A[X])
  \arrow[d]
\&
\St(\Phi,A[X])
  \arrow[d]
\\
\St(\Phi,A[X^{\pm1}])
  \arrow[r, "\chi_{\omega,X}"']
  \arrow[rr, bend right=15, "\mathrm{id}"']
\&
\St(\Phi,A[X^{\pm1}])
  \arrow[r, "\chi_{-\omega,X}"']
\&
\St(\Phi,A[X^{\pm1}])
\end{tikzcd}
\]
\end{frame}

\begin{frame}{Definition of an automorphism of $V/\sim_W$: second attempt}

Recall that $\mathrm{W}(\Phi, A) = \langle w_\alpha(u) \mid u \in A^\times \rangle \leq \St(\Phi, A)$.

\pause Set $\omega = \pm \varpi_k$. Denote by $V_\omega$ the subset of $V$ consisting of those triples $v = (x, y, z)$ for which $x \in N_\omega\cdot \mathrm{W}(\Phi, A)$.
\pause We define the function $\sigma^\pm \colon V_\omega \to V_{-\omega}$ via the following formula:
\begin{equation} \label{eq:sigma-def} (x_1 \cdot x_2, y, z) \mapsto (\delta_\omega(x_1)\cdot x_2, x_2^{-1} \cdot \chi_{\omega, X}(x_2 \cdot y),\  \widetilde{\chi}_{\omega, X}(z)), \end{equation}
where $x_1 \in N_\omega$, $x_2 \in \mathrm{W}(\Phi, A)$,\ $y \in B$, $z \in G_M$, and
$\chi_{\omega, X}$, $\widetilde{\chi}_{\omega, X}$ are weight automorphisms of $G$ and $G_M$.

\pause \textbf{Lemma 2.7:} $x_2 \cdot \chi_{\omega, X}(x_2) \in H(\Phi, A[X, X^{-1}]) = \langle h_\alpha(u) \rangle\leq B$.
\end{frame}

\begin{frame}{Definition of an automorphism of $V/\sim_W$: continued}
 $V_\omega$ is transversal wrt to $\star$: in every $\star$-orbit one can find $v \in V_\omega$.
 \pause \begin{lemma}[2.4, Bruhat decomposition]
  $\St(\Phi,A) =\UU(\Phi^\pm, A)\cdot \mathrm{W}(\Phi, A) \cdot \UU(\Phi^+, A) \cdot \overline{\St}(\Phi, A, M).$
 \end{lemma}
 \pause \begin{corollary}
  \scriptsize { $\St(\Phi,A[X]) =\UU(\Phi^\pm, A) \cdot \St(\Phi, A[X], XA[X]) \cdot \mathrm{W}(\Phi, A) \cdot \UU(\Phi^+, A) \cdot \overline{\St}(\Phi, A, M).$ }
 \end{corollary}

 \pause For every $v$ there exists $(a, b, 1) \in U^+ \times G_M^+\times B_M^+$ such that $w \star v = (x, y, z)$ and $x \in V_\omega$.

\end{frame}


\begin{frame}{Plan of the proof:}
\begin{enumerate}
 \pause \item Prove that $\delta_{\pm \varpi_k}$ exist;
 \pause \item Prove \eqref{eq:sigma-def} is well-defined and descends to $V/\sim_W$;
 \pause \item Interpret $G$ in terms of $G^+$ and $\delta_\omega$.
\end{enumerate}
\end{frame}


\begin{frame}
\begin{lemma}[4.18, easy] For $v \in V_\omega$, $v = (x_1 \cdot x_2, y, z)$ the right-hand side of\eqref{eq:sigma-def} does not depend on the decomposition $x = x_1 \cdot x_2$.
\end{lemma}
\pause \textbf{Idea:} Follows from $N_\omega \cap W(\Phi, A) = \langle w_\alpha(u) \mid \alpha \in \Delta_k, u\in A^\times\rangle$.

\pause \begin{prop}[4.21, hard] \label{prop:v-correctness2}
For $\omega = \pm \varpi_k$ and every $w \in W$, $v \in V_\omega$ such that $v' = w \star v \in V_\omega$ there exists $w' \in W$ such that
$\delta_\omega(w \star v) = w' \star \delta_\omega(v)$.
\end{prop}
\pause \textbf{Idea:} 2-pages long direct computation that relies on Chevalley--Matsumoto decomposition and explicit matrix calculations in microweight representations of $G$.
\end{frame}

\section{A presentation of $G$}

\begin{frame} \tableofcontents[currentsection] \end{frame}

\begin{frame}{Certain presentation $\St(\Phi, A[X, X^{-1}])$.}
Denote by $G'$ the amalgamated product of $\St(\Phi, A[X])$ and the infinite cyclic group $\mathbb{Z}=\langle \sigma \rangle$ amalgamated over the subgroup generated by the following relations:
\begin{align}
{}^\sigma x_{\alpha}(f) = & x_{\alpha} (Xf), & \alpha \in \Sigma^+_k, f \in A[X], \label{eq:sigma-sigma-plus} \\
x_{\beta}(f)^ \sigma     =& x_{\beta} (Xf), & \beta \in \Sigma^-_k, f \in A[X], \label{eq:sigma-sigma-minus} \\
[\sigma,\, x_\gamma(f)]   =& 1, & \gamma \in \Delta_k, f \in A[X]. \label{eq:sigma-delta}
\end{align}

\end{frame}

\begin{frame}{A presentation $\St(\Phi, A[X, X^{-1}])$.}

\begin{prop}[3.5]
For $\Phi$ and $G'$ as above there exists a group homomorphism
\[
\varphi \colon \St(\Phi,A[X,X^{-1}]) \to G'
\]
making the following diagram commute:
\[
\begin{tikzcd}[ampersand replacement=\&, column sep=2.8em, row sep=2.2em]
\& \St(\Phi,A[X,X^{-1}]) \arrow[rd, dashed, "\varphi"] \& \\
\St(\Phi,A[X]) \arrow[ru] \arrow[rr] \&\& G'
\end{tikzcd}
\]
\end{prop}
\end{frame}

\begin{frame}{Kac--Moody--Steinberg groups}
Recal that $C\in M_n(\mathbb{Z})$ is a GCM iff 

a) $c_{ii}=2$, b) $c_{ij}\leq 0$ if $i\neq j$, c) $c_{ij}=0$ iff $c_{ji}=0$, d) $c = d \cdot s$ (rational diagonal times rational symmetric). 

\pause To such $C$ one can associate a (possibly infinite) root system $\Phi = \Phi_C$ with Weyl group $W(\Phi)$. 
\pause A root $\alpha \in \Phi$ is called \textit{real} if it lies in the orbit of some simple root $\alpha_i \in \Pi$ under the action of the Weyl group $W(\Phi)$.

\pause A pair of real roots $\alpha, \beta$ is called \textit{prenilpotent} if there exists an element $w \in W(\Phi)$ such that both $w\cdot \alpha$ and $w \cdot \beta$ belong to $\Phi^+.$

\pause In particular $(\mathbb{Z}_{>0} \alpha + \mathbb{Z}_{>0}\beta)\cap \Phi \subseteq \Phi^\mathrm{re}$.
\end{frame}

\begin{frame}{Kac--Moody--Steinberg groups, continued}
$\St(C, A)$ generated by $x_\alpha(a)$, where $a \in A$ and $\alpha$ ranges over the set $\Phi^\mathrm{re}$ of real roots of $\Phi_C$.
The set of defining relations for $\St(C, A)$ consists of the following two families of relations:
\begin{align}
x_\alpha(a) x_\alpha(b) & =  x_\alpha(a + b), \label{SKM1} \\
[x_\alpha(a), x_\beta(b)] & =  \prod \limits_{\substack{i\alpha+j\beta \in \Phi^\mathrm{re} \\ i, j > 0 }} x_{i\alpha+j\beta}(N_{\alpha, \beta}^{i, j} \cdot a^i b^j), \label{SKM2}
\end{align}
where $a, b \in A$ and $\alpha, \beta$ range over all {\it prenilpotent} pairs of real roots of $\Phi$.
\end{frame}

\begin{frame}{Kac--Moody--Steinberg groups of affine type}
Let $\widetilde{\Phi}^{(1)}$ be the affine unfolded root system.
\pause \begin{lemma}
$\St(\widetilde{\Phi}^{(1)}, A) \cong \St(\Phi, A[X, X^{-1}])$.
\end{lemma}
\pause \textbf{Idea:}
$\widetilde{\Phi} \cong \Phi \times \mathbb{Z}.$ $(\alpha, n), (\beta, m)$ is prenilpotent iff $\alpha \neq - \beta$.
\begin{align*}
x_{(\alpha, m)}(a)\cdot x_{(\alpha, m)}(b)&=x_{(\alpha, m)}(a+b), \\
[x_{(\alpha, m)}(a),\,x_{(\beta, n)}(b)]  &=x_{(\alpha+\beta, n+m)}(N_{\alpha,\beta} \cdot ab),\text{ for }\alpha+\beta\in\Phi, \\
[x_{(\alpha, m)}(a),\,x_{(\beta, n)}(b)]  &=1,\text{ for }\alpha+\beta\not\in\Phi\cup0.
\end{align*}
\pause $x_{(\alpha, m)}(a) \mapsto x_\alpha(aX^m)$ is an isomorphism with the inverse given by
$x_\alpha(a_{n}X^n + \ldots + a_m X^m) \mapsto \prod_{i=n}^m x_{(\alpha, i)}(a_i)$, $a_i \in A$, $n \leq m$, $n, m\in \mathbb{Z}.$
\end{frame}

\begin{frame}{Curtis--Tits presentation, relations (simply-laced case)}
\renewcommand{\arraystretch}{1.4}

\begin{tabular}{p{0.35\textwidth} p{0.6\textwidth}}

\textbf{Condition} & \textbf{Relations} \\ \hline

all vertices $i$ on $D(C)$ &
\scriptsize{$\begin{aligned}
X_i(t)X_i(u) &= X_i(t+u) \\
[S_i^2, X_i(t)] &= 1 \\
S_i &= X_i(1)\, S_i\, X_i(1)\, S_i^{-1}\, X_i(1)
\end{aligned}$} \\
\hline
\pause
all pairs $(i,j)$ of vertices with $i \neq j$ unjoined &
\scriptsize{$\begin{aligned}
S_i S_j &= S_j S_i \\
[S_i, X_j(t)] &= 1 \\
[X_i(t), X_j(u)] &= 1
\end{aligned}$} \\
\hline
\pause
all pairs $(i,j)$ with $i \neq j$ joined by a single edge &
\scriptsize{
$\begin{aligned}
S_i S_j S_i &= S_j S_i S_j \\
S_i^2 S_j S_i^{-2} &= S_j^{-1} \\
X_i(t) S_j S_i &= S_j S_i X_i(t) \\
S_i^2 X_j(t) S_i^{-2} &= X_j(t)^{-1} \\
[X_i(t), S_i X_j(u) S_i^{-1}] &= 1 \\
[X_i(t), X_j(u)] &= S_i X_j(tu) S_i^{-1}
\end{aligned}$} \\

\end{tabular}
\end{frame}

\begin{frame}{Curtis--Tits presentation, statement}

Notice that $S_i^2 \neq 1$, instead they behave like "toric" elements. So the S-part of this presentation is something that is a quotient of the braid group (but still a nontrivial extension of the Weyl group).

\pause \begin{theorem}[Allcock, 2016]
 Assume that $C$ is $3$-spherical (or irreducible affine of rank $\geq 3$). Then $\St(C, R)$ has a presentation by means of generators $X_i(a)$, $a \in R$, $S_i$ ($i$ varies over the vertices of $D(C)$) and relations like above.
\end{theorem}
\end{frame}

\begin{frame}{Corollary 3.3}

For a finite irr, type ADE $\Phi$ and a comm. $A$ the Steinberg group $\St(\Phi, A[X, X^{-1}])$ is isomorphic to the amalgamated product of $\St(\Phi, A[X]) * \langle S \rangle$ over the subgroup generated by the following relations: 

\scriptsize 

\begin{columns}[T,onlytextwidth]
\begin{column}{0.48\textwidth}
\begin{align*}
[S^2, x_{\alpha_0}(aX)] &= 1,
\\[0.3em]
x_{\alpha_0}(X)\cdot {}^{S}x_{\alpha_0}(X)\cdot x_{\alpha_0}(X) &= S,
\\[0.3em]
[S, w_{\alpha_i}(1)] &= 1,
\\[0.3em]
[S, x_{\alpha_i}(a)] &= 1,
\\[0.3em]
S \cdot w_{\alpha_j}(1) \cdot S
&= w_{\alpha_j}(1) \cdot S \cdot w_{\alpha_j}(1),
\\[0.3em]
{}^{S^2} w_{\alpha_j}(1) &= w_{\alpha_j}(-1),
\end{align*}
\end{column}

\begin{column}{0.48\textwidth}
\begin{align*}
{}^{w_{\alpha_j}^2(1)} S &= S^{-1},
\\[0.3em]
x_{\alpha_j}(a)\cdot S \cdot w_{\alpha_j}(1)
&= S \cdot w_{\alpha_j}(1) \cdot x_{\alpha_0}(aX),
\\[0.3em]
x_{\alpha_0}(aX)\cdot w_{\alpha_j}(1)\cdot S
&= w_{\alpha_j}(1)\cdot S \cdot x_{\alpha_j}(a),
\\[0.3em]
{}^{S^2} x_{\alpha_j}(a) &= x_{\alpha_j}(-a),
\\[0.3em]
[x_{\alpha_0}(aX), {}^{S} x_{\alpha_j}(b)] &= 1,
\\[0.3em]
[x_{\alpha_0}(aX), x_{\alpha_j}(b)]
&= {}^{S} x_{\alpha_j}(ab).
\end{align*}
\end{column}
\end{columns}

$\alpha_0 := - \alpha_\mathrm{max}$, $j$ is the vertex adjacent to $0$, $i$ is a vertex unjoined with $0$, $a \in A$
\end{frame}

\begin{frame}{Sketch of the proof of Proposition 3.5}
Need to show that the following homomorphism $\varphi$ exists:
$\St(\Phi, A[X]) * \langle S \rangle / \langle\text{above page}\rangle \to \St(\Phi, A[X]) * \langle \sigma \rangle / \langle \eqref{eq:sigma-sigma-plus}\text{ -- }\eqref{eq:sigma-delta} \rangle$

\pause Set
\[\varphi(S) := x_{\alpha_0}(X) \cdot x_{-\alpha_0}(-1)^\sigma \cdot x_{\alpha_0}(X).\]

\pause Then we directly check that every relation from previous page is satisfied in $G'$.

\medskip \pause 
Example of such verification:
\scriptsize
\begin{multline*}
\varphi(x_{\alpha_0}(X) \cdot {}^{S} x_{\alpha_0}(X) \cdot x_{\alpha_0}(X)) = x_{\alpha_0}(1)^\sigma \cdot {}^{w_{\alpha_0}(1)^\sigma} x_{\alpha_0}(1)^\sigma \cdot x_{\alpha_0}(1)^\sigma = \\
= \left( x_{\alpha_0}(1) \cdot {}^{w_{\alpha_0}(1)} x_{\alpha_0}(1) \cdot x_{\alpha_0}(1)\right)^\sigma = \left(x_{\alpha_0}(1) \cdot x_{-\alpha_0}(-1) \cdot x_{\alpha_0}(1)\right)^\sigma = w_{\alpha_0}(1)^\sigma = \varphi(S).
\end{multline*}
\end{frame}

\section{The existence of $\delta_\omega$.}

\begin{frame} \tableofcontents[currentsection] \end{frame}

\begin{frame}{Ingredient 1: relative Curtis--Tits}
We want to construct $\delta_{\varpi_k} \colon N_{\varpi_k} \to N_{-\varpi_k}$.
The universal property of semidirect products reduces the problem to the construction of
$\delta_{\varpi_k}$ restricted to $\St(\Phi, A[X], XA[X])$.

\pause \begin{thm}[Lavrenov-S. 2017] \label{thm:relPres}The group $\St(\Phi, R, I)$ is isomorphic to the free product of relative Steinberg groups $\St(\Psi, R, I)$, where $\Psi \subseteq \Phi$ is of type $\rA_3$ amalgamated over embeddings of Tits generators $z_\alpha(s, \xi)$ into various subsystems $\{\alpha\} \to \Psi$. \end{thm}
\end{frame}

\begin{frame}{Ingredient 2: reduction to the case $\rA_3$}
Reduce the problem to the case $\Phi = \rA_3$, $\omega = \varpi_1$.

\pause If $\varpi_1$-pair exists for $\St(\rA_3, A[X], XA[X])$ then it exists for $\omega = \varpi_3$ and $\varpi_2$ ($\delta_{\varpi_2}$ can be expressed through $\delta_{\varpi_1}$ and $\delta_{\varpi_3}$).

\pause If $\omega$-pair exists then $w \cdot \omega$-pair exists for any $w \in W(\Phi)$:
\[ \sigma(\overline{w} \cdot \omega)(g) = w \cdot \sigma(\omega)(w^{-1} g w) \cdot w^{-1} \]

\pause Now if $\rA_3 \cong \Psi \subseteq \Phi$ then either $\Psi \subseteq \Delta_k$ (in which case $\delta_{\varpi_k}$ restricted to $\St(\Psi, A[X], XA[X])$ concides with the identity) or there exists a microweight $\omega \in P(\Psi)$ such that for every $\alpha \in \Psi$ one has $\langle \omega, \alpha \rangle = \langle \varpi_k, \alpha \rangle$.
\end{frame}

\begin{frame}{Ingredient 3: Yet another presentation}
\begin{prop}[Lavrenov--S., 2017]
\label{prop:rel-presentation}
Assume that $I$ is a splitting ideal of a commutative ring $R$.
Then for any $n\geq 4$ the group $\St(n,\,R,\,I)$ can be presented by means of two families of generators $F(u,\,v)$, $S(v,\,u)$
(where $u\in \E(n,\,R)e_1,$ and $v\in I^n$ are such that $u^{t}v=0$) subject to the following relations: $t(u, v) = e_n + u v^t$.
\begin{align*}
&F(u,\,v)\cdot F(u,\,w)=F(u,\,v+w), \\
&S(u,\,v)\cdot S(w,\,v)=S(u+w,\,v), \\
&F(u,\,v)\cdot F(u',\,v') \cdot F(u,\,v)^{-1}=F(t(u,\,v)u',\,t(v,\,u)^{-1} v'), \\
&F(me_1,\,m^{*}e_{2}a)=S(me_{1}a,\,m^{*}e_{2}),\ \text{for all $a\in I$,}\, m \in \E(n, R).
\end{align*}
\end{prop}
\end{frame}

\begin{frame}
We denote by $D(u)$ the subset of $R^n$ consisting of all columns $v$ which are orthogonal to $u$
(i.\,e. $u^{t} v = 0$) and have at least two zero entries.
\pause Recall from 3.2 of van der Kallen's paper that for every $u, v, w \in R^n$ such that $u^t v = 0$ there is a decomposition of $(w^t u) \cdot v$ into a sum of elements of $D(u)$, called \textit{canonical decomposition}:
\pause \begin{equation}
\label{eq:canonical} (w^tu) \cdot v=\sum_{i<j}u_{ij} c_{ij}(v, w),
\end{equation}
where $u_{ij} =e_i u_j-e_j u_i \in D(u)$ and $c_{ij}(v, w) =v_i w_j-v_j w_i \in R$.
\end{frame}

\begin{frame}
\begin{lemma}[van der Kallen]
\label{lem:xsmall-properties}
Let $v, w \in R^n$ s.t. $v^t w = 0$ and assume that either $v$ or $w$ has at least one zero entry.
One can define $x(v, w) \in \St(n, R)$ s.t. $\pi(x(v, w)) = t(v, w) = e + vw^t$ and
\begin{itemize}
\pause \item If $v$ or $w$ has at least $2$ zero entries, then $x(v, wa) = x(va, w)$ for $a\in R$.
\pause \item If $w_1$ and $w_2$ have at least two zero entries of which at least one entry is common
then $x(v, w_1) \cdot x(v, w_2) = x(v, w_1+w_2)$ and $x(w_1, v) \cdot x(w_2, v) = x(w_1 + w_2, v)$.
\pause \item If $v$, $v'$ are simultaneously orthogonal to $w$ and $w'$, and the elements $w, w'$ both have at least two zero entries then
$[x(v, w),\ x(v', w')] = 1$.
\pause \item If $g = x_{ij}(a)$ is a generator of the Steinberg group and $v$ or $w$ has at least two zero entries then
$g \cdot x(v, w) \cdot g^{-1} = x(gv, g^*w)$.
\end{itemize}
\end{lemma}
\end{frame}


\begin{frame}{Ingredient 4: elements $X$ and $Y$}
$R \subseteq S$ ring ext, 
$u \in \mathrm{Um}(n, R)$, $v \in I^n$ s.t. $u^{t}v = 0$.
Let $v = \sum_r v^r$ be some decomposition of $v$ into a finite sum $v^r \in D(u) \cap I^n$.
Under our assumptions on $u$ and $v$ such decomposition always exists.

\pause Let $d = \mathrm{diag}(d_1, \ldots, d_n) \in T(n, S) \leq \mathrm{GL}(n, S)$ s.t.
$d^{-1} \cdot u \in R^n,\ d \cdot I^n \subseteq R^n.$
Under these assumptions we set
\begin{equation*} X^d(u, v) :=\prod_i x(d^{-1}u, dv^r), \end{equation*}

\pause Now let $d' \in T(n, S)$ be such that $d' u\in R^n,\ {d'}^{-1} \cdot I^n \subseteq R^n$.
Under these assumptions set
\begin{equation*} Y^{d'}(v, u) := \prod_i x({d'}^{-1} v^r, {d'}u). \end{equation*}
\end{frame}

\begin{frame}{The definition of $\delta_{\varpi_1}$}
The elements $X^d(u, v)$ and $Y^{d'}(v, u)$ are well-defined, do not depend on the choice of decomposition of $v$ \pause and, moreover, for $g \in \St(n, R)$, $m = d\pi(g)d^{-1} \in \E(n, R)$ satisfy:
\begin{equation*}
g \cdot X^d(u, v) \cdot g^{-1} = X^d(mu, m^*v) \text{ and } g \cdot Y^d(v, u) \cdot g^{-1} = Y^d(mv, m^*u)
\end{equation*}

\pause Set $d_1 := \mathrm{diag}(X, 1, \ldots 1)$.
\begin{align*}
\delta_{\varpi_1} (F(u, v) ) := X^{d_1^{-1}}(u, v), & \quad \delta_{\varpi_1} (S(v, u)) := Y^{d_1^{-1}\cdot X}(v, u), \\
\delta_{-\varpi_1} (F(u, v) ) := X^{d_1 \cdot X^{-1}}(u, v),& \quad \delta_{-\varpi_1} (S(v, u)) := Y^{d_1}(v, u). 
\end{align*}
\pause It is hard to verify the $F=S$ type relation, at first we it prove only for the "big-cell" elements $m$, then use Gauss decomposition.
\end{frame}

\begin{frame}
\[ \text{Thank you for attention!} \]
\end{frame}

\end{document}
